\documentclass[quick,pro]{debate}
\motion{向孩子解释死亡}
\stance{正方}{更应该用现实}
\begin{document}
\debatetitle
\begin{thesisbox}[一句话立论]孩子来问死亡的时候问的是事实，只有先把事实说清楚，后面的温柔才有地方落脚。\end{thesisbox}
\section{定义与判准（标准表述）}
\begin{fields}
\field{现实}{以死亡的事实属性为解释的主要内容，核心三句话：不会再回来、身体停止运作、每个人都会死。语气可以温柔。
}
\field{故事}{以叙事、比喻、拟人为主要载体的解释，字面内容不要求与事实一一对应。
}
\field{孩子}{三到十二岁，重心三到八岁（提问最早三岁半，死亡概念五到七岁成形）。
}
\field{判准}{哪一种方式更能让孩子既正确地知道发生了什么，又承受得住它、并且愿意继续跟大人谈。两项合起来算，谁整体更优谁赢。
}
\field{对方提偏向定义或判准时的 30 秒口径}{对方把判准改成“孩子当下是否被安抚”：这条标准把“说的是不是真的”排除了，“爷爷出差了过几年就回来”也能得分，评委不会接受让欺骗自动过关的标准。对方把“故事”定义成“不违背事实的叙事”：那你已经在用现实了，我方欢迎。
}
\end{fields}
\section{我方论点}
\begin{longtable}{@{}L{\dimexpr0.055\linewidth-1.50\tabcolsep}L{\dimexpr0.155\linewidth-1.50\tabcolsep}L{\dimexpr0.397\linewidth-1.50\tabcolsep}L{\dimexpr0.394\linewidth-1.50\tabcolsep}@{}}
\toprule \thead{标签} & \thead{一句话} & \thead{核心数据 / 学理 / 案例} & \thead{出处}\\\midrule\endhead
孩子问的是事实 & 孩子问的是事实问题，用故事回答，他拿到的不是答案，是一个会被按字面理解的新说法 & 墨西哥 61 名 3.5–6.9 岁儿童：死亡提问 61.2\% 问死因，无一人问宗教；家长回答 22.0\% 含宗教 & Gutiérrez 等 2020，《Child Development》\\
同上 & 字面误解会造成具体伤害 & 美国儿科学会点名“睡着了”“去了很远的地方”，因孩子按字面理解，可能怕睡觉、怕出门 & HealthyChildren.org，Schonfeld 与 Nasir，2024 年 6 月 26 日更新\\
先有现实才安心 & 孩子先有生物学理解，再用它约束比喻；越早有这份理解越不怕死 & 训练实验：60 名学龄前儿童，40 人受训 20 人对照，训练组对身体与死亡的理解都显著提升 & Slaughter 与 Lyons 2003，《Cognitive Psychology》46(1): 1–30\\
同上 & 恐惧随理解下降 & 90 名四到八岁儿童，控制年龄与一般焦虑后，生物学理解越成熟死亡恐惧越低（相关，非因果） & Slaughter 与 Griffiths 2007，《Clin Child Psychol Psychiatry》12(4): 525–535\\
同上（丧亲场景） & 当年没说清楚，影响跟到成年 & 33 位童年丧亲成年人的叙事研究：当年没给清楚诚实的信息，对成年后的信任与自我价值有负面影响 & Ellis 等 2013，《J R Soc Med》106(2): 57–67\\
\bottomrule\end{longtable}
\section{对方论点 → 一句话反驳}
\begin{longtable}{@{}L{\dimexpr0.188\linewidth-1.33\tabcolsep}L{\dimexpr0.257\linewidth-1.33\tabcolsep}L{\dimexpr0.555\linewidth-1.33\tabcolsep}@{}}
\toprule \thead{对方会说} & \thead{我方一句话回} & \thead{对方追问时}\\\midrule\endhead
三到六岁接不住生物学解释 & 接不住完整的生物学，不等于接不住“不会再回来”这五个字 & 对方自己引的河南数据里，不可逆性成熟率 67.8\% 是五项最高，最能接住的被当成接不住的理由\\
故事不占用准确性（墨西哥零效应） & 你的样本上限只到六岁十一个月，代价还没到结账的时候 & 英国 136 人样本：8–11 岁宗教家庭生物学模型 23\%，非宗教 77\%\\
叙事比说明文更好懂更好记 & 那份元分析测的是读课文，不是懂死亡 & 而且好记不等于记对；Ganea 等 2014 显示拟人化绘本让孩子学到的事实更少\\
故事维系与逝者的联结 & 联结靠记忆，不靠比喻 & 《獾的礼物》的联结全靠回忆，一句谎话都没有，那正是我方的例子\\
\bottomrule\end{longtable}
\section{必问与必答}
\begin{points}{必问 （对方不答就点名，自由辩前两分钟问完）}
\item 孩子问“爷爷还会回来吗”，你方的故事怎么答这一句？（一辩稿抛出，全场追）
\item 你方的故事允许字面上与事实不符吗？
\item 你方默认的孩子几岁？八岁的孩子还用故事吗？
\item 孩子把故事当真之后要纠正，纠正用的是现实还是故事？
\end{points}
\begin{points}{必答 （15 秒内正面回）}
\item 你要给四岁孩子讲器官衰竭？ → 不是。三句话，第一句是“不会再回来”。
\item 讲现实把孩子吓哭了？ → 抱住他。哭不是失败；发现大人说的不算数才是失败。
\item 你反对绘本吗？ → 不反对。反对的是用绘本代替那三句话。
\item 《獾的礼物》算故事吗？ → 算，好故事，一句假话都没有，所以它站我方这边。
\item 训练实验只有 60 人？ → 是，对照设计，我方不夸大；对方连这样的设计都拿不出。
\end{points}
\section{自由辩战场概要}
\begin{longtable}{@{}L{\dimexpr0.135\linewidth-1.50\tabcolsep}L{\dimexpr0.330\linewidth-1.50\tabcolsep}L{\dimexpr0.162\linewidth-1.50\tabcolsep}L{\dimexpr0.373\linewidth-1.50\tabcolsep}@{}}
\toprule \thead{战场} & \thead{开场问题} & \thead{目标} & \thead{收束语}\\\midrule\endhead
一 孩子问了什么 & 孩子问“爷爷还会回来吗”，你方怎么答这一句？ & 让评委记住提问与回答的错配 & 一整场下来，对方没有正面答过这一句\\
二 故事有没有底 & 你方的故事允许字面与事实不符吗？ & 逼对方二选一 & 对方的故事要么需要现实兜底，要么会被孩子当真\\
三 谁在被安慰 & 这句“奶奶在天上看着你”，安慰的是谁？ & 为结辩的价值层铺路 & 我们争的不是要不要温柔，是温柔要不要建在真话上面\\
\bottomrule\end{longtable}
\section{如果……就……}
\begin{contingency}
\ifthen{对方改定义、改判准，或拿出没预料到的数据}{定义判准按上面的 30 秒口径打；没见过的数据先问机构、年份、样本年龄。}
\end{contingency}
\end{document}
